\documentclass[14pt,a4paper]{extarticle}
\usepackage[T2A]{fontenc}
\usepackage[utf8]{inputenc}
\usepackage[english,russian]{babel}
\usepackage{lmodern}
\usepackage{geometry}
\usepackage{array}
\usepackage{tabularx}
\usepackage{longtable}
\usepackage{ltablex}
\usepackage{booktabs}
\usepackage{hyperref}
\usepackage{indentfirst}
\usepackage{microtype}
\usepackage{graphicx}
\usepackage{setspace}
\usepackage{listings}
\usepackage{graphicx}
\usepackage{caption}
\usepackage{float}
\usepackage{tocloft}
\usepackage{titlesec}
\usepackage{newtxtext,newtxmath}

\titleformat{\section}        % форматируем \section
  {\normalfont\bfseries}      % стиль: обычный + жирный
  {\thesection.}              % номер секции с точкой
  {2mm}                       % расстояние между номером и текстом
  {}                           % перед заголовком можно добавить что-то
\titlespacing*{\section}{1.5cm}{*3}{*1}  % {отступ слева}{отступ сверху}{отступ снизу}

\captionsetup{justification=centering}
\renewcommand{\ttdefault}{cmtt}
\keepXColumns
\setlength{\parindent}{3em}
\newgeometry{left=2.5cm,right=2.5cm,top=3cm,bottom=2.5cm}

\renewcommand{\contentsname}{\hspace*{4em}{\smallСодержание}}
\setlength{\cftbeforesecskip}{0.5em}
\renewcommand{\cftsecleader}{\cftdotfill{\cftdotsep}}
\renewcommand{\cftsecaftersnum}{.\hspace{2mm}}

\title{Отчёт по культорологии}
\author{Выполнил студент группы М-РИТ-25-1, Станиславчук Сергей}

\begin{document}

\setstretch{1.5}
\noindent
\begin{minipage}[t]{0.01\textwidth}
    \includegraphics{~/programmer/++Programmer/5_1/LSTU.jpg}
\end{minipage}%
\begin{minipage}[c]{1.1\linewidth}
    \centering
    {\small
    МИНИСТЕРСТВО НАУКИ И ВЫСШЕГО ОБРАЗОВАНИЯ РФ \\[1em]
    ФЕДЕРАЛЬНОЕ ГОСУДАРСТВЕННОЕ БЮДЖЕТНОЕ \\
    ОБРАЗОВАТЕЛЬНОЕ УЧРЕЖДЕНИЕ ВЫСШЕГО ОБРАЗОВАНИЯ \\
    «ЛИПЕЦКИЙ ГОСУДАРСТВЕННЫЙ ТЕХНИЧЕСКИЙ УНИВЕРСИТЕТ»
    }
\end{minipage}%

\vspace{3em}

\noindent
\hspace{1em}%
\begin{minipage}[t]{0.8\textwidth}
\textbf{Институт} \hspace{2em} компьютерных наук \\
\textbf{Кафедра} \hspace{2em} автоматизированных систем управления
\end{minipage}


\vspace{6em}

\begin{center}
\textbf{ЗАДАНИЕ 8} \\
по дисциплине \\ <<Межкультурное взаимодействие>>
\end{center}

\vspace{6em}

\noindent
\begin{tabular*}{\textwidth}{@{\extracolsep{\fill}} l c r }
    Студент М-РИТ-25-1 & \makebox[4cm]{\hrulefill} & Станиславчук С.М. \\[2em]
Руководитель & \\
Ст. преподаватель & \makebox[4cm]{\hrulefill} & Болдырихин О.В. \\
\end{tabular*}

\vfill
\begin{center}
    Липецк, 2025
\end{center}

\thispagestyle{empty}

\newpage
\tableofcontents
\newpage

\section{Задание 1. Национальные модели лидерства}
\begin{sloppypar}
Модель лидерства напрямую зависит от культурных установок общества, его
    отношения к власти, ответственности и социальной иерархии. В разных странах
    представления о том, каким должен быть лидер, формируются историческими
    традициями, экономическим устройством и национальными ценностями.

    В странах с низкой дистанцией власти, например в Скандинавии, лидера
    рассматривают как равного члену команды. Он действует через убеждение,
    сотрудничество и коллективные решения. Здесь ценится горизонтальная
    структура, открытость, доверие и минимальная формальность.

    В США или Великобритании лидерство строится на личной инициативе,
    независимости, способности брать на себя ответственность и достигать
    результата. Лидер — это мотиватор, который вдохновляет и направляет, но не
    контролирует каждый шаг.

    В странах с высокой дистанцией власти, таких как Китай, Индия или
    большинство арабских государств, лидер — это фигура авторитетная, часто
    неприкасаемая. Он принимает ключевые решения, ожидает уважения и
    безусловного признания статуса. Команда ориентируется на выполнение
    указаний и соблюдение иерархии.

    В Японии лидерство коллективистское: руководитель отвечает за гармонию и
    стабильность, решения принимаются медленно, после долгих консультаций.
    Лидер — координатор, который обеспечивает баланс интересов, а не
    доминирует.
\end{sloppypar}

\section{Задание 2. Представления о статусе и лидерстве в различных культурых.}
\begin{sloppypar}
Статус и лидерство интерпретируются по-разному в зависимости от культурной
    среды. В одних обществах статус определяется достижениями и качествами
    личности, в других — происхождением, возрастом или социальным положением.

    В культурах, ориентированных на достижения (США, Канада, Германия), статус
    формируется за счёт компетентности, профессионализма, инициативности. Лидер
    — это человек, который доказал эффективность и готов брать на себя
    ответственность. Здесь важна демонстрация навыков, открытость и готовность
    аргументировать решения.

    В культурах, ориентированных на предписанный статус (Япония, Китай, Индия),
    большую роль играют возраст, опыт, социальное происхождение и занимаемая
    должность. Лидер определяется не столько личными качествами, сколько
    соответствием традиционным нормам. Уважение идёт «сверху вниз» и редко
    ставится под сомнение.

    В индивидуалистических обществах лидер — самостоятельная фигура, которая
    принимает решения самостоятельно и несёт за них личную ответственность. В
    коллективистских культурах лидер рассматривается как представитель группы,
    а его задача — сохранять единство, гармонию и внутреннюю стабильность.
\end{sloppypar}

\section{Задание 3. Факторы, воздействующие на эффективность лидерства в различных культурах.}
\begin{sloppypar}
Эффективность лидерства во многом определяется культурным контекстом, в котором
действует руководитель. Первый ключевой фактор — уровень дистанции власти. В
странах с высокой дистанцией власти подчинённые ожидают жёсткой вертикали,
чёткого контроля и однозначных указаний. Лидер, пытающийся руководить
демократично и неформально, будет восприниматься как слабый. В культурах с
низкой дистанцией власти, наоборот, авторитарный стиль вызывает сопротивление и
недоверие.

Второй фактор — ориентация на индивидуализм или коллективизм. В
индивидуалистических культурах лидер эффективен тогда, когда стимулирует личные
достижения, поощряет инициативу и создаёт возможности для самореализации. В
коллективистских культурах ценится лидер, который укрепляет группу,
распределяет ответственность, поддерживает сплочённость и избегает открытых
конфликтов.

Сильное влияние оказывает и стиль коммуникации. Там, где предпочитают прямой
стиль, лидер обязан говорить ясно и однозначно. В культурах непрямой
коммуникации эффективность строится на умении читать контекст, уважать намёки и
соблюдать ритуалы вежливости. Неспособность правильно интерпретировать эти
сигналы резко снижает авторитет.

Последний важный фактор — отношение к неопределённости. В странах с высокой
нетерпимостью к неопределённости лидер должен обеспечивать предсказуемость,
структурированность и правила. В странах с низкой нетерпимостью к
неопределённости допускаются гибкость, эксперименты и более свободный подход к
управлению.
\end{sloppypar}

\section{Задание 4. Классификация проблем межгосударственного различия в управлении человеческими ресурсами}
\begin{sloppypar}
Различия между государствами создают целый набор проблем при управлении
    персоналом в международных компаниях. Первая группа проблем связана с
    нормативно-правовыми различиями. Страны отличаются трудовым
    законодательством, правилами найма и увольнения, требованиями к охране
    труда, социальными гарантиями. Это заставляет компании адаптировать
    кадровую политику для каждой юрисдикции отдельно.

    Вторая группа проблем обусловлена культурными различиями. Уровень
    иерархичности, способы принятия решений, предпочтительные формы мотивации,
    отношение к критике, инициативности и ответственности — всё это формирует
    разные стандарты поведения работников. Непонимание этих норм порождает
    конфликты, снижает мотивацию и эффективность команд.

    Третья группа — различия в ожиданиях работников. В одних странах сотрудники
    ориентированы на стабильность, долгосрочную занятость и социальную
    защищённость. В других — на быстрый карьерный рост, гибкость и высокую
    мобильность. Универсальная система мотивации в таких условиях плохо
    работает.

    Четвёртая группа проблем — языковые барьеры и особенности коммуникации.
    Даже при знании общего языка различаются стиль общения, уровень прямоты,
    способы выражения несогласия и формы рабочего взаимодействия. Это влияет на
    качество командной работы и точность передачи информации.

    Пятая группа — различия в управленческих моделях. В одних государствах
    доминирует директивное управление, в других — командное или гибридное.
    Международная компания вынуждена балансировать между корпоративным
    стандартом и локальными ожиданиями, что далеко не всегда просто.
\end{sloppypar}


\end{document}
